\chapter{The take-home}
\label{ch:takehome}

\section*{What this chapter is for}

\reported{A take-home appears in about a third of loops that disclose their
process, typically at two to three hours, and usually in place of an on-site
round rather than in addition to
one.}{field-guide-process} Nobody watches you do
it, which removes every signal the live round was generating and replaces it
with what you choose to leave behind.

\section{Start with evaluation}

The single strongest piece of guidance available about these assessments:
\reported{not starting with evaluation is reported as a red flag, and analysed
signals name starting with evals before the main logic as the first
thing.}{field-guide-take-home}

That is worth reading twice, because it inverts the instinct. The instinct is
to build the thing and then, if time remains, to check it. The reported
expectation is the reverse: decide what correct means, encode it, then build
towards it.

\begin{kitmethod}
\textbf{First half-hour, before the main logic:}

\textbf{1.} Write down what correct output looks like for five or six inputs
\dash{} by hand, from the brief. These are your golden cases.

\textbf{2.} Write the thing that checks them. It will fail. That is the point:
you have a target rather than an opinion.

\textbf{3.} Include at least one case that must \emph{not} happen \dash{} a
tool that must not be called, a field that must not appear, an action that must
not be taken without confirmation. A check that only looks for the right answer
misses half of what can go wrong.

\textbf{4.} Then build.
\end{kitmethod}

\judgement{This costs half an hour of a three-hour budget and it changes what
the whole submission communicates. A repository whose first commit is a failing
evaluation says something about the author that no amount of polish in the
final one does.}

\section{The other signals}

\reported{The analysed evaluation signals for these assessments are: modular
structure and error handling; production awareness \dash{} caching, monitoring,
cost, security; edge cases handled even where tests were not requested; and
documented decisions.}{field-guide-take-home}

Three of those four are things you would do anyway. The fourth is the one that
needs a deliberate decision, because nobody asks for it.

\textbf{Documented decisions} means a short file \dash{} four or five entries
\dash{} each naming a decision, the alternative, and why. Not a design
document. A list of forks with the reason for the turn taken.

And, from Chapter~\ref{ch:ledger}: the register of what you did not do. In a
take-home it is more valuable than anywhere else, because the reader's default
assumption about anything missing is that you did not think of it.

\section{Scope, when nobody is stopping you}

The stated budget is the budget. Working eight hours on a three-hour assessment
is a worse submission, not a better one: it fails the honesty test if you do not
say so, and if you do say so it reports that you cannot scope yourself.

\begin{kitmethod}
\textbf{Timebox it in writing, in the repository.} ``Three hours, spent
roughly: 30 minutes evaluation, 90 minutes the main path, 30 minutes the second
case, 30 minutes documentation.'' A reader now knows what they are looking at.

\textbf{Cut breadth, never completeness.} Chapter~\ref{ch:slices}, unchanged.

\textbf{List what the next two hours would have gone on.} It costs four lines
and answers the question the reader is forming anyway.
\end{kitmethod}

\section{The defence round is the assessment}

\reported{The defence conversation afterwards is reported to matter more than
the code alone.}{field-guide-take-home}

So the submission is written to be defended. Before you send it, prepare
answers to five questions \dash{} they are asked nearly every time:

\begin{itemize}
  \item Why this structure rather than the obvious alternative?
  \item What is the weakest part?
  \item What would break first at a hundred times the volume?
  \item What did you not have time for, and what would you do first?
  \item How do you know it works?
\end{itemize}

\judgement{The fifth is the one that separates submissions, and it is why the
first half-hour went where it did. ``Here is the check, here is what it covers,
here is what it does not'' is an answer. ``I tested it manually'' is the
absence of one.}

\section{A video walkthrough}

\reported{Submitting with a short video walkthrough is among the reported
hygiene practices for these assessments.}{field-guide-get-hired}

Three or four minutes: what it does, the one interesting decision, and what you
left out. \judgement{It recovers most of the signal the live round would have
produced \dash{} how you talk about your own work, what you choose to point at,
whether you volunteer the weakness \dash{} which is the whole thing the format
gave up by not watching you.}

\begin{kittrap}
The trap specific to this format is polish. Nobody is watching, the clock is
soft, and it is easy to spend the last hour on presentation \dash{} formatting,
a nicer interface, a longer README.

That hour buys almost nothing. The same hour spent on one more golden case, one
more entry in the decisions file, or one honest paragraph about what is not
covered changes what the submission is evidence \emph{of}.
\end{kittrap}

\begin{drill}
Take any past exercise and redo its first thirty minutes under this
instruction: write the golden cases and the failing check, and nothing else.
Notice how much the brief's ambiguity resists being turned into concrete
expected outputs \dash{} that resistance is the discovery work of
Chapter~\ref{ch:discovery} arriving in a format where you cannot ask anybody.
\end{drill}

\section*{What this chapter does not claim}

The reported signals come from an analysis of take-home repositories and the
feedback attached to them, in one corpus. That is a real body of evidence about
a real set of assessments, and it is not a survey of what every employer scores.

The recommendation to spend the first sixth of the budget on evaluation is a
judgement built on that reporting. The reporting says starting with evaluation
matters; the fraction is the author's.
